\chapter{Consigna 2: Respuesta en Frecuencia y Compensación}

\section{Ítem 1: Diagrama de Bode de la Planta sin Compensar}

De acuerdo con los parámetros obligatorios asignados al Grupo 5 en la Tabla 1 de las consignas, la planta a caracterizar
posee una ganancia estática $K = 900$ y tres polos reales simples cuyas frecuencias de corte corresponden a $f_1 = 12\,\si{\kilo\hertz}$,
$f_2 = 70\,\si{\kilo\hertz}$ y $f_3 = 140\,\si{\kilo\hertz}$. La función de transferencia en el dominio de Laplace adopta
la siguiente expresión canónica:

\begin{equation}
    G(s) = \frac{K}{\left(1 + \frac{s}{2\pi f_1}\right) \left(1 + \frac{s}{2\pi f_2}\right) \left(1 + \frac{s}{2\pi f_3}\right)}
    \label{eq.ej2:planta_simbolica}
\end{equation}

Sustituyendo los valores asignados, resulta:

\begin{equation}
    G(s) = \frac{900}{\left(1 + \frac{s}{2\pi \cdot 12 \times 10^3}\right) \left(1 + \frac{s}{2\pi \cdot 70 \times 10^3}\right) \left(1 + \frac{s}{2\pi \cdot 140 \times 10^3}\right)}
    \label{eq.ej2:planta_numerica}
\end{equation}

En la Figura \ref{fig.ej2:bode_planta} se presenta el diagrama de Bode de la planta con su respuesta en frecuencia exacta, el cual constituye la base gráfica para las determinaciones de estabilidad requeridas.

\begin{figure}[!ht]
    \centering
    \begin{tikzpicture}
    \pgfplotsset{
        bode plot/.style={
            /pgf/number format/read comma as period,
            grid=both,
            major grid style={solid, gray!45},
            minor grid style={dotted, gray!35},
            axis line style={black, line width=0.6pt},
            axis x line* = bottom,
            axis y line* = left,
            extra description/.code={
                \draw[gray!45] (rel axis cs:0,1) -- (rel axis cs:1,1) -- (rel axis cs:1,0);
            },
            width=15.5cm,
            height=6.5cm,
            xmin=100, xmax=10000000,
            enlarge x limits=false,
            tick label style={font=\small},
            label style={font=\normalsize},
        }
    }

    % Subplot 1: Magnitud
    \begin{semilogxaxis}[
        bode plot,
        ylabel={Magnitud [\si{\decibel}]},
        ymin=-70, ymax=75,
        ytick={-60, -40, -20, 0, 20, 40, 60},
        xticklabel=\empty,
        legend pos=south west,
        legend style={font=\small, fill=white, fill opacity=0.9, text opacity=1},
        name=magplot,
    ]
        % Curva exacta
        \addplot[line width=1.3pt, blue!85!black] table [col sep=comma, x=frequency, y=mag_exact] {sources/plots/bode_planta.csv};
        \addlegendentry{$G(s)$}

        % Nivel 0 dB
        \addplot[black!70, densely dotted, line width=0.8pt] coordinates {(100, 0) (10000000, 0)};
        \node[black!70, anchor=south east, font=\scriptsize] at (axis cs:9000000, 1.5) {$0\,\text{dB}$};

        % Ancho de banda a -3 dB
        \draw[teal!80!black, dashed, line width=1.0pt] (axis cs:11593, -70) -- (axis cs:11593, 56.07);
        \addplot[only marks, mark=*, mark size=2.5pt, teal!80!black] coordinates {(11593, 56.07)};
        \node[teal!80!black, fill=white, fill opacity=0.9, text opacity=1, inner sep=2pt, font=\small, anchor=south] 
            at (axis cs:11593, 62) {$BW = 11,59\,\text{kHz}$};
        \draw[teal!80!black, -stealth, line width=0.8pt] (axis cs:11593, 61.5) -- (axis cs:11593, 56.7);

        % Guía vertical f_45 (MF = 45°)
        \draw[purple, dashed, line width=1.0pt] (axis cs:52901, -70) -- (axis cs:52901, 43.44);
        \addplot[only marks, mark=*, mark size=2.6pt, purple] coordinates {(52901, 43.44)};
        \node[purple, fill=white, fill opacity=0.9, text opacity=1, inner sep=1.5pt, font=\small, anchor=south] 
            at (axis cs:52901, 45.5) {$43,44\,\text{dB}$};

        % Guía vertical f_180
        \draw[magenta!80!black, dashed, line width=1.0pt] (axis cs:110993, -70) -- (axis cs:110993, 32.14);
        \addplot[only marks, mark=square*, mark size=2.4pt, magenta!80!black] coordinates {(110993, 32.14)};
        \node[magenta!80!black, fill=white, fill opacity=0.9, text opacity=1, inner sep=1.5pt, font=\small, anchor=north east] 
            at (axis cs:105000, 30.5) {$32,14\,\text{dB}$};

        % Cota de Margen de Ganancia para MF = 45°
        \draw[orange!90!black, dotted, line width=0.8pt] (axis cs:52901, 43.44) -- (axis cs:230000, 43.44);
        \draw[orange!90!black, dotted, line width=0.8pt] (axis cs:110993, 32.14) -- (axis cs:230000, 32.14);
        \draw[stealth-stealth, orange!90!black, line width=1.2pt] 
            (axis cs:200000, 32.14) -- (axis cs:200000, 43.44)
            node[midway, right=2pt, font=\small, fill=white, fill opacity=0.9, text opacity=1, inner sep=1.5pt] {$MG_{45} = 11,30\,\text{dB}$};
    \end{semilogxaxis}

    % Subplot 2: Fase (posicionado estrictamente DEBAJO de magplot)
    \begin{semilogxaxis}[
        bode plot,
        xlabel={$f$ [\si{\hertz}]},
        ylabel={Fase [$^\circ$]},
        ymin=-280, ymax=10,
        ytick={-270, -225, -180, -135, -90, -45, 0},
        at={(magplot.south west)},
        anchor=north west,
        yshift=-0.7cm,
    ]
        % Curva exacta
        \addplot[line width=1.3pt, blue!85!black] table [col sep=comma, x=frequency, y=phase_exact] {sources/plots/bode_planta.csv};

        % Línea de referencia -135° (MF = 45°)
        \addplot[purple, dashed, line width=1.0pt] coordinates {(100, -135) (10000000, -135)};
        \node[purple, fill=white, fill opacity=0.85, text opacity=1, anchor=south east, font=\small, inner sep=1.5pt] 
            at (axis cs:9000000, -132) {$-135^\circ\;(MF=45^\circ)$};

        % Línea de referencia -180°
        \addplot[red!80!black, dotted, line width=1.0pt] coordinates {(100, -180) (10000000, -180)};
        \node[red!80!black, fill=white, fill opacity=0.85, text opacity=1, anchor=south east, font=\small, inner sep=1.5pt] 
            at (axis cs:9000000, -177) {$-180^\circ$};

        % Cruce f_45
        \draw[purple, dashed, line width=1.0pt] (axis cs:52901, -280) -- (axis cs:52901, -135);
        \addplot[only marks, mark=*, mark size=2.6pt, purple] coordinates {(52901, -135)};
        \node[purple, fill=white, fill opacity=0.9, text opacity=1, inner sep=2pt, font=\small, anchor=north] 
            at (axis cs:52901, -143) {$f_{45} = 52,90\,\text{kHz}$};

        % Cruce f_180
        \draw[magenta!80!black, dashed, line width=1.0pt] (axis cs:110993, -280) -- (axis cs:110993, -180);
        \addplot[only marks, mark=square*, mark size=2.4pt, magenta!80!black] coordinates {(110993, -180)};
        \node[magenta!80!black, fill=white, fill opacity=0.9, text opacity=1, inner sep=2pt, font=\small, anchor=north] 
            at (axis cs:110993, -188) {$f_{180} = 110,99\,\text{kHz}$};
    \end{semilogxaxis}
\end{tikzpicture}

    \caption{Diagrama de Bode de magnitud y fase de la planta sin compensar $G(s)$, con indicación gráfica del ancho de banda ($BW$), frecuencia para $MF = 45^\circ$ ($f_{45}$) y margen de ganancia ($MG_{45}$).}
    \label{fig.ej2:bode_planta}
\end{figure}

\section{Ítem 2: Obtención Gráfica de Frecuencia de Corte y Margen de Ganancia para MF = 45°}

A partir de la inspección directa del diagrama de Bode de la Figura \ref{fig.ej2:bode_planta}, se determinan de forma
gráfica los parámetros dinámicos y de estabilidad del sistema:

\begin{enumerate}
    \item \textbf{Ancho de banda a lazo abierto ($BW = f_{-3\,\si{\decibel}}$):} En el gráfico de magnitud, la asíntota de baja
    frecuencia se sitúa en $59,08\,\si{\decibel}$. Descendiendo una cota vertical de $3\,\si{\decibel}$ hasta el nivel de
    $56,08\,\si{\decibel}$, se traza una línea horizontal que corta a la respuesta real de la planta. Proyectando dicho punto de
    intersección verticalmente hacia el eje horizontal de frecuencias, se lee en el gráfico:
    \begin{equation}
        BW = f_{-3\,\si{\decibel}} \approx 11,59\,\si{\kilo\hertz}
    \end{equation}
    Este valor evidencia que el ancho de banda a lazo abierto se encuentra casi exclusivamente gobernado por el primer polo
    ($f_1 = 12\,\si{\kilo\hertz}$).
    \item \textbf{Margen de ganancia para $MF = 45^\circ$ ($MG_{45}$):} Sobre la curva de fase, la condición de margen de
    fase $MF = 45^\circ$ se alcanza cuando el desfasaje vale $-135^\circ$ ($\phi(f_{45}) = -180^\circ + 45^\circ = -135^\circ$),
    lo que proyectado sobre el eje de frecuencias sitúa $f_{45} \approx 52,90\,\si{\kilo\hertz}$. Sobre esa misma vertical
    en el gráfico de magnitud, la ganancia de la planta alcanza:
    \begin{equation}
        |G(f_{45})|_{\si{\decibel}} \approx 43,44\,\si{\decibel}
    \end{equation}
    Asimismo, se ubica la intersección de la fase con el nivel de $-180^\circ$ (cruce de fase), leyendo sobre el eje horizontal
    la cota $f_{180} \approx 110,99\,\si{\kilo\hertz}$, donde la ganancia de la planta vale:
    \begin{equation}
        |G(f_{180})|_{\si{\decibel}} \approx 32,14\,\si{\decibel}
    \end{equation}
    El margen de ganancia disponible cuando el lazo se calibra para entregar $MF = 45^\circ$ se obtiene gráficamente como la
    distancia vertical en decibeles entre el nivel de lazo en $f_{45}$ y el nivel en $f_{180}$, acotada mediante la flecha
    naranja de la Figura \ref{fig.ej2:bode_planta}:
    \begin{equation}
        MG_{45} = |G(f_{45})|_{\si{\decibel}} - |G(f_{180})|_{\si{\decibel}} = 43,44\,\si{\decibel} - 32,14\,\si{\decibel} = 11,30\,\si{\decibel}
    \end{equation}
    Nótese que el margen de ganancia directo de la planta no atenuada respecto del nivel de $0\,\si{\decibel}$ es negativo
    ($MG = -32,14\,\si{\decibel}$), lo que ratifica que la planta a lazo cerrado con realimentación unitaria es incondicionalmente
    inestable.
\end{enumerate}

\clearpage

\section{Ítem 3: Compensador de Retardo}

El compensador de retardo asignado al Grupo 5 se proyecta mediante un filtro pasa bajos pasivo RC constituido por una
resistencia $R$ en serie y un capacitor $C$ en derivación a masa, de acuerdo con el esquema circuital de la Figura
\ref{crkt.ej2:retardo}.

\begin{figure}[!ht]
    \centering
    \begin{circuitikz}[american, scale=1.0]
    \draw (0,2) node[left] {$V_i$} 
          to[short, o-] (0.5,2)
          to[R, l=$R$, -*] (3,2)
          to[short, -o] (4.2,2) node[right] {$V_o$};
    \draw (3,2) 
          to[C, l=$C$] (3,0) node[ground] {};
    \draw (0,0) node[left] {} 
          to[short, o-] (0.5,0) node[ground] {};
    \draw (4.2,0) node[right] {} 
          to[short, o-] (3.7,0) node[ground] {};
\end{circuitikz}

    \caption{Esquema circuital del compensador de retardo (filtro pasa bajos RC).}
    \label{crkt.ej2:retardo}
\end{figure}

Aplicando divisor de tensión en el dominio de Laplace sobre el capacitor $C$, la función de transferencia de la red resulta:

\begin{equation}
    G_c(s) = \frac{V_o(s)}{V_i(s)} = \frac{\frac{1}{sC}}{R + \frac{1}{sC}} = \frac{1}{1 + sRC} = \frac{1}{1 + \frac{s}{2\pi f_d}}
    \label{eq.ej2:fdt_comp_retardo}
\end{equation}

donde la frecuencia característica asignada al Grupo 5 es $f_d = 1,2\,\si{\kilo\hertz}$. Fijando la capacidad en $C = 100\,\si{\nano\farad} = 10^{-7}\,\si{\farad}$,
se despeja el valor de la resistencia $R$:

\begin{equation}
    R = \frac{1}{2\pi f_d C} = \frac{1}{2\pi \cdot 1,2 \times 10^3\,\si{\hertz} \cdot 100 \times 10^{-9}\,\si{\farad}} \approx 1326,29\,\si{\ohm}
    \label{eq.ej2:calculo_R_retardo}
\end{equation}

El valor comercial normalizado más próximo al $1\%$ de tolerancia (serie E96) es $R = 1,33\,\si{\kilo\ohm}$ ($1,30\,\si{\kilo\ohm}$
en serie E24 al $5\%$).

Conectando el compensador en cascada con la planta, la función de transferencia total resulta:

\begin{equation}
    G_{\text{ret}}(s) = G_c(s) \cdot G(s) = \frac{900}{\left(1 + \frac{s}{2\pi \cdot 1,2 \times 10^3}\right) \left(1 + \frac{s}{2\pi \cdot 12 \times 10^3}\right) \left(1 + \frac{s}{2\pi \cdot 70 \times 10^3}\right) \left(1 + \frac{s}{2\pi \cdot 140 \times 10^3}\right)}
    \label{eq.ej2:fdt_sistema_retardo}
\end{equation}

En la Figura \ref{fig.ej2:bode_retardo} se ilustra el diagrama de Bode resultante.

\begin{figure}[!ht]
    \centering
    \begin{tikzpicture}
    \pgfplotsset{
        bode plot/.style={
            /pgf/number format/read comma as period,
            grid=both,
            major grid style={solid, gray!45},
            minor grid style={dotted, gray!35},
            axis line style={black, line width=0.6pt},
            axis x line* = bottom,
            axis y line* = left,
            extra description/.code={
                \draw[gray!45] (rel axis cs:0,1) -- (rel axis cs:1,1) -- (rel axis cs:1,0);
            },
            width=15.5cm,
            height=6.5cm,
            xmin=100, xmax=10000000,
            enlarge x limits=false,
            tick label style={font=\small},
            label style={font=\normalsize},
        }
    }

    % Subplot 1: Magnitud
    \begin{semilogxaxis}[
        bode plot,
        ylabel={Magnitud [\si{\decibel}]},
        ymin=-100, ymax=75,
        ytick={-80, -60, -40, -20, 0, 20, 40, 60},
        xticklabel=\empty,
        legend pos=south west,
        legend style={font=\small, fill=white, fill opacity=0.9, text opacity=1},
        name=magplot,
    ]
        % Curva exacta
        \addplot[line width=1.3pt, blue!85!black] table [col sep=comma, x=frequency, y=mag_exact] {sources/plots/bode_retardo.csv};
        \addlegendentry{$G_{\text{ret}}(s)$}

        % Nivel 0 dB
        \addplot[black!70, densely dotted, line width=0.8pt] coordinates {(100, 0) (10000000, 0)};
        \node[black!70, anchor=south east, font=\scriptsize] at (axis cs:9000000, 1.5) {$0\,\text{dB}$};

        % Ancho de banda a -3 dB
        \draw[teal!80!black, dashed, line width=1.0pt] (axis cs:1188, -100) -- (axis cs:1188, 56.07);
        \addplot[only marks, mark=*, mark size=2.5pt, teal!80!black] coordinates {(1188, 56.07)};
        \node[teal!80!black, fill=white, fill opacity=0.9, text opacity=1, inner sep=2pt, font=\small, anchor=south] 
            at (axis cs:1188, 62) {$BW = 1,19\,\text{kHz}$};
        \draw[teal!80!black, -stealth, line width=0.8pt] (axis cs:1188, 61.5) -- (axis cs:1188, 56.7);

        % Guía vertical f_45
        \draw[purple, dashed, line width=1.0pt] (axis cs:9962, -100) -- (axis cs:9962, 38.25);
        \addplot[only marks, mark=*, mark size=2.6pt, purple] coordinates {(9962, 38.25)};
        \node[purple, fill=white, fill opacity=0.9, text opacity=1, inner sep=1.5pt, font=\small, anchor=south] 
            at (axis cs:9962, 40.5) {$38,25\,\text{dB}$};

        % Guía vertical f_180
        \draw[magenta!80!black, dashed, line width=1.0pt] (axis cs:24354, -100) -- (axis cs:24354, 25.21);
        \addplot[only marks, mark=square*, mark size=2.4pt, magenta!80!black] coordinates {(24354, 25.21)};
        \node[magenta!80!black, fill=white, fill opacity=0.9, text opacity=1, inner sep=1.5pt, font=\small, anchor=north east] 
            at (axis cs:23000, 23.5) {$25,21\,\text{dB}$};

        % Cota de Margen de Ganancia para MF = 45°
        \draw[orange!90!black, dotted, line width=0.8pt] (axis cs:9962, 38.25) -- (axis cs:55000, 38.25);
        \draw[orange!90!black, dotted, line width=0.8pt] (axis cs:24354, 25.21) -- (axis cs:55000, 25.21);
        \draw[stealth-stealth, orange!90!black, line width=1.2pt] 
            (axis cs:50000, 25.21) -- (axis cs:50000, 38.25)
            node[midway, right=2pt, font=\small, fill=white, fill opacity=0.9, text opacity=1, inner sep=1.5pt] {$MG_{45} = 13,04\,\text{dB}$};
    \end{semilogxaxis}

    % Subplot 2: Fase (posicionado estrictamente DEBAJO de magplot)
    \begin{semilogxaxis}[
        bode plot,
        xlabel={$f$ [\si{\hertz}]},
        ylabel={Fase [$^\circ$]},
        ymin=-370, ymax=10,
        ytick={-360, -315, -270, -225, -180, -135, -90, -45, 0},
        at={(magplot.south west)},
        anchor=north west,
        yshift=-0.7cm,
    ]
        % Curva exacta
        \addplot[line width=1.3pt, blue!85!black] table [col sep=comma, x=frequency, y=phase_exact] {sources/plots/bode_retardo.csv};

        % Línea de referencia -135° (MF = 45°)
        \addplot[purple, dashed, line width=1.0pt] coordinates {(100, -135) (10000000, -135)};
        \node[purple, fill=white, fill opacity=0.85, text opacity=1, anchor=south east, font=\small, inner sep=1.5pt] 
            at (axis cs:9000000, -132) {$-135^\circ\;(MF=45^\circ)$};

        % Línea de referencia -180°
        \addplot[red!80!black, dotted, line width=1.0pt] coordinates {(100, -180) (10000000, -180)};
        \node[red!80!black, fill=white, fill opacity=0.85, text opacity=1, anchor=south east, font=\small, inner sep=1.5pt] 
            at (axis cs:9000000, -177) {$-180^\circ$};

        % Cruce f_45
        \draw[purple, dashed, line width=1.0pt] (axis cs:9962, -370) -- (axis cs:9962, -135);
        \addplot[only marks, mark=*, mark size=2.6pt, purple] coordinates {(9962, -135)};
        \node[purple, fill=white, fill opacity=0.9, text opacity=1, inner sep=2pt, font=\small, anchor=north] 
            at (axis cs:9962, -143) {$f_{45} = 9,96\,\text{kHz}$};

        % Cruce f_180
        \draw[magenta!80!black, dashed, line width=1.0pt] (axis cs:24354, -370) -- (axis cs:24354, -180);
        \addplot[only marks, mark=square*, mark size=2.4pt, magenta!80!black] coordinates {(24354, -180)};
        \node[magenta!80!black, fill=white, fill opacity=0.9, text opacity=1, inner sep=2pt, font=\small, anchor=north] 
            at (axis cs:24354, -188) {$f_{180} = 24,35\,\text{kHz}$};
    \end{semilogxaxis}
\end{tikzpicture}

    \caption{Diagrama de Bode del sistema compensado con red de retardo $G_{\text{ret}}(s)$, indicando gráficamente el ancho de banda ($BW$), la frecuencia para $MF = 45^\circ$ ($f_{45}$) y el margen de ganancia ($MG_{45}$).}
    \label{fig.ej2:bode_retardo}
\end{figure}

A partir de la lectura gráfica sobre la Figura \ref{fig.ej2:bode_retardo}, se extraen los siguientes parámetros:

\begin{enumerate}
    \item \textbf{Ancho de banda ($BW$):} La caída de $3\,\si{\decibel}$ respecto del nivel de continua ($59,08\,\si{\decibel} - 3\,\si{\decibel} = 56,08\,\si{\decibel}$)
    se ubica ahora en la frecuencia del nuevo polo dominante introducido por el compensador:
    \begin{equation}
        BW = f_{-3\,\si{\decibel}} \approx 1,19\,\si{\kilo\hertz}
    \end{equation}
    lo que evidencia una reducción drástica del ancho de banda útil a lazo abierto a menos del $10\%$ del valor original.

    \item \textbf{Margen de ganancia para $MF = 45^\circ$ ($MG_{45}$):} Sobre la curva de fase, el punto donde la fase
    alcanza $-135^\circ$ (correspondiente a $MF = 45^\circ$) se sitúa en $f_{45} \approx 9,96\,\si{\kilo\hertz}$, donde la curva
    de magnitud entrega:
    \begin{equation}
        |G_{\text{ret}}(f_{45})|_{\si{\decibel}} \approx 38,25\,\si{\decibel}
    \end{equation}
    Por su parte, la traza de fase corta a $-180^\circ$ en la frecuencia gráfica:
    \begin{equation}
        f_{180} \approx 24,35\,\si{\kilo\hertz}
    \end{equation}
    donde la magnitud del sistema compensado vale:
    \begin{equation}
        |G_{\text{ret}}(f_{180})|_{\si{\decibel}} \approx 25,21\,\si{\decibel}
    \end{equation}
    El margen de ganancia para $MF = 45^\circ$ se determina como la diferencia vertical entre ambas cotas:
    \begin{equation}
        MG_{45} = |G_{\text{ret}}(f_{45})|_{\si{\decibel}} - |G_{\text{ret}}(f_{180})|_{\si{\decibel}} = 38,25\,\si{\decibel} - 25,21\,\si{\decibel} = 13,04\,\si{\decibel}
    \end{equation}
    El margen de ganancia directo respecto de $0\,\si{\decibel}$ se ubica en $-25,21\,\si{\decibel}$.
\end{enumerate}

\clearpage

\section{Ítem 4: Compensador de Atraso-Adelanto}

La red de compensación de atraso-adelanto se diseña a partir de la topología serie conformada por los elementos $R_1 \to R_2 \to C$,
tomando la salida de tensión en el nodo intermedio ubicado antes de $R_2$, tal como se representa en el esquemático de la
Figura \ref{crkt.ej2:atraso_adelanto}.

\begin{figure}[!ht]
    \centering
    \begin{circuitikz}[american, scale=1.0]
    \draw (0,3) node[left] {$V_i$} 
          to[short, o-] (0.5,3)
          to[R, l=$R_1$, -*] (3,3)
          to[short, -o] (4.5,3) node[right] {$V_o$};
    \draw (3,3) 
          to[R, l=$R_2$, *-] (3,1.5)
          to[C, l=$C$] (3,0) node[ground] {};
    \draw (0,0) node[left] {} 
          to[short, o-] (0.5,0) node[ground] {};
    \draw (4.5,0) node[right] {} 
          to[short, o-] (4.0,0) node[ground] {};
\end{circuitikz}

    \caption{Esquema circuital del compensador de atraso-adelanto (red serie $R_1 \to R_2 \to C$).}
    \label{crkt.ej2:atraso_adelanto}
\end{figure}

La impedancia de la rama en derivación formada por $R_2$ y $C$ en serie responde a:

\begin{equation}
    Z_2(s) = R_2 + \frac{1}{sC} = \frac{sR_2C + 1}{sC}
\end{equation}

mientras que la impedancia de entrada en circuito abierto vale:

\begin{equation}
    Z_{\text{in}}(s) = R_1 + Z_2(s) = R_1 + R_2 + \frac{1}{sC} = \frac{s(R_1 + R_2)C + 1}{sC}
\end{equation}

Aplicando divisor de impedancias, la función de transferencia resulta:

\begin{equation}
    G_{c2}(s) = \frac{V_o(s)}{V_i(s)} = \frac{Z_2(s)}{Z_{\text{in}}(s)} = \frac{sR_2C + 1}{s(R_1 + R_2)C + 1} = \frac{1 + \frac{s}{2\pi f_z}}{1 + \frac{s}{2\pi f_p}}
    \label{eq.ej2:fdt_atraso_adelanto_canonica}
\end{equation}

de donde se desprenden las expresiones de las frecuencias singulares:

\begin{equation}
    f_z = \frac{1}{2\pi R_2 C}, \qquad f_p = \frac{1}{2\pi (R_1 + R_2) C}
    \label{eq.ej2:fz_fp_expresiones}
\end{equation}

Dado que $R_1 > 0$, se comprueba rigurosamente que $f_p < f_z$. Para el Grupo 5, los valores asignados son $f_p = 220\,\si{\kilo\hertz}$
y $f_z = 1,1\,\si{\mega\hertz} = 1100\,\si{\kilo\hertz}$. Fijando $C = 100\,\si{\pico\farad} = 10^{-10}\,\si{\farad}$, se procede al
cálculo secuencial de las resistencias:

\begin{enumerate}
    \item \textbf{Cálculo de $R_2$ (a partir de $f_z$):}
    \begin{equation}
        R_2 = \frac{1}{2\pi f_z C} = \frac{1}{2\pi \cdot 1,1 \times 10^6\,\si{\hertz} \cdot 100 \times 10^{-12}\,\si{\farad}} \approx 1446,86\,\si{\ohm}
        \label{eq.ej2:calculo_R2}
    \end{equation}
    (Valor normalizado comercial al $1\%$: $R_2 = 1,45\,\si{\kilo\ohm}$).

    \item \textbf{Cálculo de $R_1$ (a partir de $f_p$ y $R_2$):}
    \begin{equation}
        R_1 + R_2 = \frac{1}{2\pi f_p C} = \frac{1}{2\pi \cdot 220 \times 10^3\,\si{\hertz} \cdot 100 \times 10^{-12}\,\si{\farad}} \approx 7234,32\,\si{\ohm}
    \end{equation}
    Despejando $R_1$:
    \begin{equation}
        R_1 = (R_1 + R_2) - R_2 = 7234,32\,\si{\ohm} - 1446,86\,\si{\ohm} \approx 5787,45\,\si{\ohm}
        \label{eq.ej2:calculo_R1}
    \end{equation}
    (Valor normalizado comercial al $1\%$: $R_1 = 5,76\,\si{\kilo\ohm}$).
\end{enumerate}


La función de transferencia en cascada con la planta resulta:

\begin{equation}
    G_{\text{aa}}(s) = \frac{900 \left(1 + \frac{s}{2\pi \cdot 1,1 \times 10^6}\right)}{\left(1 + \frac{s}{2\pi \cdot 220 \times 10^3}\right) \left(1 + \frac{s}{2\pi \cdot 12 \times 10^3}\right) \left(1 + \frac{s}{2\pi \cdot 70 \times 10^3}\right) \left(1 + \frac{s}{2\pi \cdot 140 \times 10^3}\right)}
    \label{eq.ej2:fdt_sistema_aa}
\end{equation}

En la Figura \ref{fig.ej2:bode_aa} se exhibe el diagrama de Bode correspondiente a esta topología.

\begin{figure}[!ht]
    \centering
    \begin{tikzpicture}
    \pgfplotsset{
        bode plot/.style={
            /pgf/number format/read comma as period,
            grid=both,
            major grid style={solid, gray!45},
            minor grid style={dotted, gray!35},
            axis line style={black, line width=0.6pt},
            axis x line* = bottom,
            axis y line* = left,
            extra description/.code={
                \draw[gray!45] (rel axis cs:0,1) -- (rel axis cs:1,1) -- (rel axis cs:1,0);
            },
            width=15.5cm,
            height=6.5cm,
            xmin=100, xmax=10000000,
            enlarge x limits=false,
            tick label style={font=\small},
            label style={font=\normalsize},
        }
    }

    % Subplot 1: Magnitud
    \begin{semilogxaxis}[
        bode plot,
        ylabel={Magnitud [\si{\decibel}]},
        ymin=-70, ymax=75,
        ytick={-60, -40, -20, 0, 20, 40, 60},
        xticklabel=\empty,
        legend pos=south west,
        legend style={font=\small, fill=white, fill opacity=0.9, text opacity=1},
        name=magplot,
    ]
        % Curva exacta
        \addplot[line width=1.3pt, blue!85!black] table [col sep=comma, x=frequency, y=mag_exact] {sources/plots/bode_atraso_adelanto.csv};
        \addlegendentry{$G_{\text{aa}}(s)$}

        % Nivel 0 dB
        \addplot[black!70, densely dotted, line width=0.8pt] coordinates {(100, 0) (10000000, 0)};
        \node[black!70, anchor=south east, font=\scriptsize] at (axis cs:9000000, 1.5) {$0\,\text{dB}$};

        % Ancho de banda a -3 dB
        \draw[teal!80!black, dashed, line width=1.0pt] (axis cs:11563, -70) -- (axis cs:11563, 56.07);
        \addplot[only marks, mark=*, mark size=2.5pt, teal!80!black] coordinates {(11563, 56.07)};
        \node[teal!80!black, fill=white, fill opacity=0.9, text opacity=1, inner sep=2pt, font=\small, anchor=south] 
            at (axis cs:11563, 62) {$BW = 11,56\,\text{kHz}$};
        \draw[teal!80!black, -stealth, line width=0.8pt] (axis cs:11563, 61.5) -- (axis cs:11563, 56.7);

        % Guía vertical f_45
        \draw[purple, dashed, line width=1.0pt] (axis cs:45109, -70) -- (axis cs:45109, 45.18);
        \addplot[only marks, mark=*, mark size=2.6pt, purple] coordinates {(45109, 45.18)};
        \node[purple, fill=white, fill opacity=0.9, text opacity=1, inner sep=1.5pt, font=\small, anchor=south] 
            at (axis cs:45109, 47) {$45,18\,\text{dB}$};

        % Guía vertical f_180
        \draw[magenta!80!black, dashed, line width=1.0pt] (axis cs:84486, -70) -- (axis cs:84486, 36.22);
        \addplot[only marks, mark=square*, mark size=2.4pt, magenta!80!black] coordinates {(84486, 36.22)};
        \node[magenta!80!black, fill=white, fill opacity=0.9, text opacity=1, inner sep=1.5pt, font=\small, anchor=north east] 
            at (axis cs:80000, 34.5) {$36,22\,\text{dB}$};

        % Cota de Margen de Ganancia para MF = 45°
        \draw[orange!90!black, dotted, line width=0.8pt] (axis cs:45109, 45.18) -- (axis cs:200000, 45.18);
        \draw[orange!90!black, dotted, line width=0.8pt] (axis cs:84486, 36.22) -- (axis cs:200000, 36.22);
        \draw[stealth-stealth, orange!90!black, line width=1.2pt] 
            (axis cs:170000, 36.22) -- (axis cs:170000, 45.18)
            node[midway, right=2pt, font=\small, fill=white, fill opacity=0.9, text opacity=1, inner sep=1.5pt] {$MG_{45} = 8,96\,\text{dB}$};
    \end{semilogxaxis}

    % Subplot 2: Fase (posicionado estrictamente DEBAJO de magplot)
    \begin{semilogxaxis}[
        bode plot,
        xlabel={$f$ [\si{\hertz}]},
        ylabel={Fase [$^\circ$]},
        ymin=-280, ymax=10,
        ytick={-270, -225, -180, -135, -90, -45, 0},
        at={(magplot.south west)},
        anchor=north west,
        yshift=-0.7cm,
    ]
        % Curva exacta
        \addplot[line width=1.3pt, blue!85!black] table [col sep=comma, x=frequency, y=phase_exact] {sources/plots/bode_atraso_adelanto.csv};

        % Línea de referencia -135° (MF = 45°)
        \addplot[purple, dashed, line width=1.0pt] coordinates {(100, -135) (10000000, -135)};
        \node[purple, fill=white, fill opacity=0.85, text opacity=1, anchor=south east, font=\small, inner sep=1.5pt] 
            at (axis cs:9000000, -132) {$-135^\circ\;(MF=45^\circ)$};

        % Línea de referencia -180°
        \addplot[red!80!black, dotted, line width=1.0pt] coordinates {(100, -180) (10000000, -180)};
        \node[red!80!black, fill=white, fill opacity=0.85, text opacity=1, anchor=south east, font=\small, inner sep=1.5pt] 
            at (axis cs:9000000, -177) {$-180^\circ$};

        % Cruce f_45
        \draw[purple, dashed, line width=1.0pt] (axis cs:45109, -280) -- (axis cs:45109, -135);
        \addplot[only marks, mark=*, mark size=2.6pt, purple] coordinates {(45109, -135)};
        \node[purple, fill=white, fill opacity=0.9, text opacity=1, inner sep=2pt, font=\small, anchor=north] 
            at (axis cs:45109, -143) {$f_{45} = 45,11\,\text{kHz}$};

        % Cruce f_180
        \draw[magenta!80!black, dashed, line width=1.0pt] (axis cs:84486, -280) -- (axis cs:84486, -180);
        \addplot[only marks, mark=square*, mark size=2.4pt, magenta!80!black] coordinates {(84486, -180)};
        \node[magenta!80!black, fill=white, fill opacity=0.9, text opacity=1, inner sep=2pt, font=\small, anchor=north] 
            at (axis cs:84486, -188) {$f_{180} = 84,49\,\text{kHz}$};
    \end{semilogxaxis}
\end{tikzpicture}

    \caption{Diagrama de Bode del sistema con compensador de atraso-adelanto $G_{\text{aa}}(s)$, señalando el ancho de banda ($BW$), la frecuencia para $MF = 45^\circ$ ($f_{45}$) y el margen de ganancia ($MG_{45}$).}
    \label{fig.ej2:bode_aa}
\end{figure}

A partir del examen de la Figura \ref{fig.ej2:bode_aa}, se determinan gráficamente los siguientes valores:

\begin{enumerate}
    \item \textbf{Frecuencia y ganancia para $MF = 45^\circ$:} Localizando el desfasaje de $-135^\circ$ sobre la curva de fase,
    se proyecta hacia el eje de frecuencias, leyendo:
    \begin{equation}
        f_{45} \approx 45,11\,\si{\kilo\hertz}
    \end{equation}
    Sobre la misma vertical en el gráfico de magnitud, la ganancia correspondiente se lee directamente como:
    \begin{equation}
        |G_{\text{aa}}(f_{45})|_{\si{\decibel}} \approx 45,18\,\si{\decibel} \quad (\approx 181,50\text{ veces})
    \end{equation}
    Por consiguiente, para fijar $MF = 45^\circ$ se requiere una atenuación de lazo de $\beta = 1/181,5 \approx 0,00551$ ($-45,18\,\si{\decibel}$).

    \item \textbf{Margen de ganancia ($MG_{45}$):} El cruce de fase por $-180^\circ$ se localiza gráficamente en:
    \begin{equation}
        f_{180} \approx 84,49\,\si{\kilo\hertz}
    \end{equation}
    donde la curva de magnitud entrega:
    \begin{equation}
        |G_{\text{aa}}(f_{180})|_{\si{\decibel}} \approx 36,22\,\si{\decibel}
    \end{equation}
    El margen de ganancia para $MF = 45^\circ$ se obtiene de la cota vertical entre $f_{45}$ y $f_{180}$:
    \begin{equation}
        MG_{45} = 45,18\,\si{\decibel} - 36,22\,\si{\decibel} = 8,96\,\si{\decibel}
    \end{equation}
    mientras que el margen de ganancia directo respecto de $0\,\si{\decibel}$ se ubica en $-36,22\,\si{\decibel}$.

    \item \textbf{Ancho de banda a lazo abierto ($BW$):} La cota de $-3\,\si{\decibel}$ respecto de continua se mantiene en:
    \begin{equation}
        BW = f_{-3\,\si{\decibel}} \approx 11,56\,\si{\kilo\hertz}
    \end{equation}
    evidenciando que este compensador no degrada el ancho de banda base de la planta.
\end{enumerate}

\clearpage

\section{Ítem 5: Análisis Comparativo y Explicación de Diferencias}

La Figura \ref{fig.ej2:bode_comparativa} ilustra la superposición directa de los diagramas de Bode para la planta original
y las dos configuraciones compensadas. Por su parte, la Tabla \ref{tab.ej2:resumen_comparativo} resume las magnitudes
cuantitativas deducidas gráficamente.

\begin{figure}[!ht]
    \centering
    \begin{tikzpicture}
    \pgfplotsset{
        bode plot/.style={
            /pgf/number format/read comma as period,
            grid=both,
            major grid style={solid, gray!45},
            minor grid style={dotted, gray!35},
            axis line style={black, line width=0.6pt},
            axis x line* = bottom,
            axis y line* = left,
            extra description/.code={
                \draw[gray!45] (rel axis cs:0,1) -- (rel axis cs:1,1) -- (rel axis cs:1,0);
            },
            width=15.5cm,
            height=6.5cm,
            xmin=100, xmax=10000000,
            enlarge x limits=false,
            tick label style={font=\small},
            label style={font=\normalsize},
        }
    }

    % Subplot 1: Magnitud comparada
    \begin{semilogxaxis}[
        bode plot,
        ylabel={Magnitud [\si{\decibel}]},
        ymin=-100, ymax=75,
        ytick={-80, -60, -40, -20, 0, 20, 40, 60},
        xticklabel=\empty,
        legend pos=south west,
        legend style={font=\small, fill=white, fill opacity=0.9, text opacity=1},
        name=magplot,
    ]
        % Planta
        \addplot[line width=1.3pt, blue!85!black] table [col sep=comma, x=frequency, y=mag_planta] {sources/plots/bode_comparativa.csv};
        \addlegendentry{Planta original $G(s)$}

        % Retardo
        \addplot[line width=1.3pt, red!85!black] table [col sep=comma, x=frequency, y=mag_retardo] {sources/plots/bode_comparativa.csv};
        \addlegendentry{Comp. Retardo $G_{\text{ret}}(s)$}

        % Atraso-Adelanto
        \addplot[line width=1.3pt, teal!80!black] table [col sep=comma, x=frequency, y=mag_aa] {sources/plots/bode_comparativa.csv};
        \addlegendentry{Comp. Atraso-Adelanto $G_{\text{aa}}(s)$}

        % Nivel 0 dB
        \addplot[black!70, densely dotted, line width=0.8pt] coordinates {(100, 0) (10000000, 0)};
        \node[black!70, anchor=south east, font=\scriptsize] at (axis cs:9000000, 1.5) {$0\,\text{dB}$};

        % Proyecciones verticales f_45
        \draw[blue!85!black, dashed, line width=0.9pt] (axis cs:52901, -100) -- (axis cs:52901, 43.44);
        \draw[red!85!black, dashed, line width=0.9pt] (axis cs:9962, -100) -- (axis cs:9962, 38.25);
        \draw[teal!80!black, dashed, line width=0.9pt] (axis cs:45109, -100) -- (axis cs:45109, 45.18);

        % Marcadores f_45
        \addplot[only marks, mark=*, mark size=2.5pt, blue!85!black] coordinates {(52901, 43.44)};
        \node[blue!85!black, fill=white, fill opacity=0.9, text opacity=1, inner sep=1.5pt, font=\scriptsize, anchor=south west] 
            at (axis cs:56000, 44.5) {$43,44\,\text{dB}$};

        \addplot[only marks, mark=*, mark size=2.5pt, red!85!black] coordinates {(9962, 38.25)};
        \node[red!85!black, fill=white, fill opacity=0.9, text opacity=1, inner sep=1.5pt, font=\scriptsize, anchor=north] 
            at (axis cs:9962, 36) {$38,25\,\text{dB}$};

        \addplot[only marks, mark=*, mark size=2.5pt, teal!80!black] coordinates {(45109, 45.18)};
        \node[teal!80!black, fill=white, fill opacity=0.9, text opacity=1, inner sep=1.5pt, font=\scriptsize, anchor=south east] 
            at (axis cs:42000, 46.5) {$45,18\,\text{dB}$};
    \end{semilogxaxis}

    % Subplot 2: Fase comparada (posicionado estrictamente DEBAJO de magplot)
    \begin{semilogxaxis}[
        bode plot,
        xlabel={$f$ [\si{\hertz}]},
        ylabel={Fase [$^\circ$]},
        ymin=-370, ymax=10,
        ytick={-360, -315, -270, -225, -180, -135, -90, -45, 0},
        at={(magplot.south west)},
        anchor=north west,
        yshift=-0.7cm,
        legend pos=south west,
        legend style={font=\small, fill=white, fill opacity=0.9, text opacity=1},
    ]
        % Planta
        \addplot[line width=1.3pt, blue!85!black] table [col sep=comma, x=frequency, y=phase_planta] {sources/plots/bode_comparativa.csv};
        \addlegendentry{Planta original}

        % Retardo
        \addplot[line width=1.3pt, red!85!black] table [col sep=comma, x=frequency, y=phase_retardo] {sources/plots/bode_comparativa.csv};
        \addlegendentry{Retardo}

        % Atraso-Adelanto
        \addplot[line width=1.3pt, teal!80!black] table [col sep=comma, x=frequency, y=phase_aa] {sources/plots/bode_comparativa.csv};
        \addlegendentry{Atraso-Adelanto}

        % Referencia -135°
        \addplot[purple, dashed, line width=1.0pt] coordinates {(100, -135) (10000000, -135)};
        \node[purple, fill=white, fill opacity=0.85, text opacity=1, anchor=south east, font=\small, inner sep=1.5pt] 
            at (axis cs:9000000, -132) {$-135^\circ\;(MF=45^\circ)$};

        % Referencia -180°
        \addplot[red!80!black, dotted, line width=1.0pt] coordinates {(100, -180) (10000000, -180)};
        \node[red!80!black, fill=white, fill opacity=0.85, text opacity=1, anchor=south east, font=\small, inner sep=1.5pt] 
            at (axis cs:9000000, -177) {$-180^\circ$};

        % Proyecciones verticales f_45
        \draw[blue!85!black, dashed, line width=0.9pt] (axis cs:52901, -370) -- (axis cs:52901, -135);
        \draw[red!85!black, dashed, line width=0.9pt] (axis cs:9962, -370) -- (axis cs:9962, -135);
        \draw[teal!80!black, dashed, line width=0.9pt] (axis cs:45109, -370) -- (axis cs:45109, -135);

        % Cruces en f_45
        \addplot[only marks, mark=*, mark size=2.5pt, blue!85!black] coordinates {(52901, -135)};
        \node[blue!85!black, fill=white, fill opacity=0.85, text opacity=1, inner sep=1pt, font=\scriptsize, anchor=south west] 
            at (axis cs:56000, -130) {$f_{45}=52,90\,\text{kHz}$};

        \addplot[only marks, mark=*, mark size=2.5pt, red!85!black] coordinates {(9962, -135)};
        \node[red!85!black, fill=white, fill opacity=0.85, text opacity=1, inner sep=1pt, font=\scriptsize, anchor=north] 
            at (axis cs:9962, -143) {$f_{45}=9,96\,\text{kHz}$};

        \addplot[only marks, mark=*, mark size=2.5pt, teal!80!black] coordinates {(45109, -135)};
        \node[teal!80!black, fill=white, fill opacity=0.85, text opacity=1, inner sep=1pt, font=\scriptsize, anchor=north east] 
            at (axis cs:42000, -143) {$f_{45}=45,11\,\text{kHz}$};
    \end{semilogxaxis}
\end{tikzpicture}

    \caption{Comparativa gráfica de Bode (magnitud y fase) entre la planta original $G(s)$, con compensador de retardo $G_{\text{ret}}(s)$ y con compensador de atraso-adelanto $G_{\text{aa}}(s)$.}
    \label{fig.ej2:bode_comparativa}
\end{figure}

\begin{table}[!ht]
    \centering
    \renewcommand{\arraystretch}{0.92}
    \resizebox{\textwidth}{!}{%
    \begin{tabular}{|l|c|c|c|}
        \hline
        \textbf{Parámetro / Especificación} & \textbf{Planta Original} & \textbf{Comp. Retardo} & \textbf{Comp. Atraso-Adelanto} \\
        \hline\hline
        Ganancia estática $K$ & $900$ ($59,08\,\si{\decibel}$) & $900$ ($59,08\,\si{\decibel}$) & $900$ ($59,08\,\si{\decibel}$) \\
        Ancho de banda a $-3\,\si{\decibel}$ ($BW_{\text{LA}}$) & $11,59\,\si{\kilo\hertz}$ & $1,19\,\si{\kilo\hertz}$ & $11,56\,\si{\kilo\hertz}$ \\
        Frecuencia para $MF = 45^\circ$ ($f_{45}$) & $52,90\,\si{\kilo\hertz}$ & $9,96\,\si{\kilo\hertz}$ & $45,11\,\si{\kilo\hertz}$ \\
        Ganancia disponible en $f_{45}$ ($|G(f_{45})|$) & $43,44\,\si{\decibel}$ ($148,6\times$) & $38,25\,\si{\decibel}$ ($81,8\times$) & $45,18\,\si{\decibel}$ ($181,5\times$) \\
        Frecuencia de cruce de fase $-180^\circ$ ($f_{180}$) & $110,99\,\si{\kilo\hertz}$ & $24,35\,\si{\kilo\hertz}$ & $84,49\,\si{\kilo\hertz}$ \\
        Ganancia en cruce $-180^\circ$ ($|G(f_{180})|$) & $32,14\,\si{\decibel}$ & $25,21\,\si{\decibel}$ & $36,22\,\si{\decibel}$ \\
        Margen de Ganancia para $MF = 45^\circ$ ($MG_{45}$) & $11,30\,\si{\decibel}$ & $13,04\,\si{\decibel}$ & $8,96\,\si{\decibel}$ \\
        Margen de Ganancia directo ($MG$) & $-32,14\,\si{\decibel}$ & $-25,21\,\si{\decibel}$ & $-36,22\,\si{\decibel}$ \\
        Frecuencia de cruce $0\,\si{\decibel}$ sin atenuar ($f_{0\,\si{\decibel}}$) & $464,42\,\si{\kilo\hertz}$ & $83,94\,\si{\kilo\hertz}$ & $373,66\,\si{\kilo\hertz}$ \\
        Margen de Fase sin atenuar ($MF_{\text{directo}}$) & $-63,17^\circ$ & $-72,17^\circ$ & $-97,76^\circ$ \\
        \hline
    \end{tabular}%
    }
    \caption{Cuadro comparativo de las características obtenidas a partir de los diagramas de Bode para el Grupo 5.}
    \label{tab.ej2:resumen_comparativo}
\end{table}

A partir de la confrontación de los resultados de los Ítems 3 y 4, se explican las siguientes diferencias fundamentales:

\begin{enumerate}
    \item \textbf{Impacto sobre el ancho de banda y la frecuencia de corte:} La frecuencia de corte de la función de
    transferencia viene determinada por el punto donde la ganancia cae $3\,\si{\decibel}$ respecto del valor máximo en
    baja frecuencia ($BW = f_{-3\,\si{\decibel}}$). En la planta original sin compensar, dicha frecuencia es
    $BW = 11,59\,\si{\kilo\hertz}$, impuesta fundamentalmente por su primer polo ($f_1 = 12\,\si{\kilo\hertz}$).

    Al incorporar el compensador de retardo, el polo introducido a baja frecuencia ($f_d = 1,2\,\si{\kilo\hertz}$) actúa como
    dominante y fuerza una caída prematura de $-20\,\si{\decibel\per\decade}$, contrayendo drásticamente el ancho de banda a
    $BW = 1,19\,\si{\kilo\hertz}$ (pérdida del $89,7\%$). En marcado contraste, la red de atraso-adelanto ubica su polo en
    $f_p = 220\,\si{\kilo\hertz}$ y su cero en $f_z = 1,1\,\si{\mega\hertz}$, ambos por encima de los polos de la planta,
    preservando la frecuencia de corte prácticamente intacta en $BW = 11,56\,\si{\kilo\hertz}$ y conservando plenamente la
    velocidad de respuesta de la planta.

    Por otro lado, la condición de margen de fase $MF = 45^\circ$ (desfasaje de $-135^\circ$) se alcanza a $f_{45} = 9,96\,\si{\kilo\hertz}$
    en el compensador de retardo y a $f_{45} = 45,11\,\si{\kilo\hertz}$ en atraso-adelanto, corroborando que este último admite
    operar a frecuencias sensiblemente mayores bajo el mismo margen de estabilidad.

    \item \textbf{Margen de estabilidad y atenuación de alta frecuencia:} El compensador de retardo entrega un margen de ganancia
    superior ($MG_{45} = 13,04\,\si{\decibel}$ frente a $8,96\,\si{\decibel}$).

\end{enumerate}
